\documentclass{beamer}
\usetheme{Madrid}
\usecolortheme{default}
\usepackage[utf8]{inputenc}
\usepackage{amsmath}
\usepackage{amsfonts}
\usepackage{amssymb}
\usepackage{tikz}

\title{CRISC Domain D - 100 More Impossible Hard Questions}
\subtitle{Information Technology and Security - Expert Level Part 2}
\author{CRISC Mastery Series}
\date{\today}

\begin{document}
	
	\begin{frame}
		\titlepage
	\end{frame}
	
	\begin{frame}{Warning}
		\begin{alertblock}{Extreme Difficulty}
			Domain D is technical but the exam tests judgment, not trivia.\\
			Every option will look plausible.\\
			Think like a risk practitioner advising on technology, not a technician.
		\end{alertblock}
	\end{frame}
	
	\begin{frame}{Outline}
		\tableofcontents
	\end{frame}
	
	% ============================================================
	\section{1. Enterprise Architecture}
	% ============================================================
	
	\begin{frame}{Q101: Architecture Governance Board}
		\begin{block}{Question}
			An architecture review board rejects a project because it violates a principle. The project sponsor appeals. What is the BEST governance path?
		\end{block}
		\begin{enumerate}[(A)]
			\item The sponsor proceeds
			\item The board decision stands unless escalated through defined exception governance
			\item The board is dissolved
			\item The vendor decides
		\end{enumerate}
	\end{frame}
	
	\begin{frame}{Q101: Explanation}
		\begin{block}{Correct Answer: B}
			\textbf{Exceptions follow governance:}
			\begin{itemize}
				\item Boards enforce principles
				\item Exceptions require escalation and approval
				\item Proceeding undermines the board
			\end{itemize}
		\end{block}
		\begin{alertblock}{Critical Trap}
			Architecture exceptions need governance, not bypass.
		\end{alertblock}
	\end{frame}
	
	\begin{frame}{Q102: TOGAF ADM Purpose}
		\begin{block}{Question}
			What is the purpose of the TOGAF Architecture Development Method?
		\end{block}
		\begin{enumerate}[(A)]
			\item Marketing plan
			\item Structured method for developing and managing architecture
			\item Network monitoring
			\item HR planning
		\end{enumerate}
	\end{frame}
	
	\begin{frame}{Q102: Explanation}
		\begin{block}{Correct Answer: B}
			\textbf{ADM structures architecture work:}
			\begin{itemize}
				\item Provides phases from vision to maintenance
				\item Iterative and adaptable
				\item Aligns architecture to strategy
			\end{itemize}
		\end{block}
		\begin{alertblock}{Critical Trap}
			ADM is the method, not a framework replacement.
		\end{alertblock}
	\end{frame}
	
	\begin{frame}{Q103: Architecture Views}
		\begin{block}{Question}
			Why does EA use multiple views?
		\end{block}
		\begin{enumerate}[(A)]
			\item To confuse stakeholders
			\item To communicate to different audiences (business, data, application, technology)
			\item To increase costs
			\item To hide issues
		\end{enumerate}
	\end{frame}
	
	\begin{frame}{Q103: Explanation}
		\begin{block}{Correct Answer: B}
			\textbf{Views match audiences:}
			\begin{itemize}
				\item Different views for different concerns
				\item Business view supports executives
				\item Technology view supports engineers
			\end{itemize}
		\end{block}
		\begin{alertblock}{Critical Trap}
			Views exist to communicate, not to obscure.
		\end{alertblock}
	\end{frame}
	
	\begin{frame}{Q104: EA Metrics}
		\begin{block}{Question}
			Which metric BEST reflects EA value?
		\end{block}
		\begin{enumerate}[(A)]
			\item Number of diagrams produced
			\item Reduction in duplication and improved alignment with strategy
			\item Meeting attendance
			\item Vendor count
		\end{enumerate}
	\end{frame}
	
	\begin{frame}{Q104: Explanation}
		\begin{block}{Correct Answer: B}
			\textbf{EA value is alignment and efficiency:}
			\begin{itemize}
				\item Reduced duplication saves cost
				\item Alignment protects value
				\item Diagrams are outputs, not outcomes
			\end{itemize}
		\end{block}
		\begin{alertblock}{Critical Trap}
			Diagrams are not the value; outcomes are.
		\end{alertblock}
	\end{frame}
	
	\begin{frame}{Q105: Reference Architectures}
		\begin{block}{Question}
			What is the MAIN benefit of a reference architecture?
		\end{block}
		\begin{enumerate}[(A)]
			\item Avoids standards
			\item Provides a reusable pattern for consistent designs
			\item Faster meetings
			\item Larger budgets
		\end{enumerate}
	\end{frame}
	
	\begin{frame}{Q105: Explanation}
		\begin{block}{Correct Answer: B}
			\textbf{Reference architectures reuse patterns:}
			\begin{itemize}
				\item Consistent design accelerates delivery
				\item Reduces rework
				\item Improves quality
			\end{itemize}
		\end{block}
		\begin{alertblock}{Critical Trap}
			Reference architectures are patterns, not mandates.
		\end{alertblock}
	\end{frame}
	
	\begin{frame}{Q106: Architecture Principle Example}
		\begin{block}{Question}
			Which is a valid architecture principle?
		\end{block}
		\begin{enumerate}[(A)]
			\item All vendors must be fired
			\item Data is a shared enterprise asset with controlled access
			\item Only one server is allowed
			\item Ignore all audits
		\end{enumerate}
	\end{frame}
	
	\begin{frame}{Q106: Explanation}
		\begin{block}{Correct Answer: B}
			\textbf{Principles must guide decisions:}
			\begin{itemize}
				\item Data as shared asset guides access design
				\item Principles are directional, not operational rules
				\item Extreme statements are not principles
			\end{itemize}
		\end{block}
		\begin{alertblock}{Critical Trap}
			Principles guide, not ban.
		\end{alertblock}
	\end{frame}
	
	\begin{frame}{Q107: Architecture Repository}
		\begin{block}{Question}
			What is the MAIN purpose of an architecture repository?
		\end{block}
		\begin{enumerate}[(A)]
			\item Store backups
			\item Centralize architecture assets, standards, and decisions
			\item Track HR records
			\item Manage payroll
		\end{enumerate}
	\end{frame}
	
	\begin{frame}{Q107: Explanation}
		\begin{block}{Correct Answer: B}
			\textbf{Repository preserves architecture knowledge:}
			\begin{itemize}
				\item Central access to assets
				\item Supports reuse and governance
				\item Prevents lost knowledge
			\end{itemize}
		\end{block}
		\begin{alertblock}{Critical Trap}
			Repository is a knowledge asset, not a backup.
		\end{alertblock}
	\end{frame}
	
	\begin{frame}{Q108: EA and Business Strategy}
		\begin{block}{Question}
			Which of the following signals that EA is aligned to business strategy?
		\end{block}
		\begin{enumerate}[(A)]
			\item IT chooses tools without business input
			\item Roadmaps reflect business objectives and priorities
			\item Only vendors influence architecture
			\item Architecture is independent of strategy
		\end{enumerate}
	\end{frame}
	
	\begin{frame}{Q108: Explanation}
		\begin{block}{Correct Answer: B}
			\textbf{Alignment shows in roadmaps:}
			\begin{itemize}
				\item Roadmaps translate strategy into action
				\item Business objectives drive technology
				\item Isolation signals misalignment
			\end{itemize}
		\end{block}
		\begin{alertblock}{Critical Trap}
			Alignment is proven by roadmaps, not documents.
		\end{alertblock}
	\end{frame}
	
	\begin{frame}{Q109: Architecture Change Management}
		\begin{block}{Question}
			What does architecture change management do?
		\end{block}
		\begin{enumerate}[(A)]
			\item Ignores changes
			\item Governs changes to the architecture and maintains alignment
			\item Backs up files
			\item Manages emails
		\end{enumerate}
	\end{frame}
	
	\begin{frame}{Q109: Explanation}
		\begin{block}{Correct Answer: B}
			\textbf{Change keeps architecture aligned:}
			\begin{itemize}
				\item Governs architecture change
				\item Prevents drift
				\item Maintains strategy alignment
			\end{itemize}
		\end{block}
		\begin{alertblock}{Critical Trap}
			Architecture changes need governance.
		\end{alertblock}
	\end{frame}
	
	\begin{frame}{Q110: EA Maturity}
		\begin{block}{Question}
			What does high EA maturity usually indicate?
		\end{block}
		\begin{enumerate}[(A)]
			\item More drawings
			\item Consistent, strategy-aligned technology decisions across the enterprise
			\item Larger budgets
			\item Fewer audits
		\end{enumerate}
	\end{frame}
	
	\begin{frame}{Q110: Explanation}
		\begin{block}{Correct Answer: B}
			\textbf{High maturity is alignment at scale:}
			\begin{itemize}
				\item Consistent decisions
				\item Strategy-driven technology
				\item Reduced duplication
			\end{itemize}
		\end{block}
		\begin{alertblock}{Critical Trap}
			Maturity is about consistency, not documents.
		\end{alertblock}
	\end{frame}
	
	% ============================================================
	\section{2. IT Operations and Change Management}
	% ============================================================
	
	\begin{frame}{Q111: Change Management KPIs}
		\begin{block}{Question}
			Which KPI BEST measures change management effectiveness?
		\end{block}
		\begin{enumerate}[(A)]
			\item Number of meetings
			\item Percentage of changes with successful outcomes and no unplanned downtime
			\item Number of tickets closed
			\item Budget spent
		\end{enumerate}
	\end{frame}
	
	\begin{frame}{Q111: Explanation}
		\begin{block}{Correct Answer: B}
			\textbf{Change success rate is the KPI:}
			\begin{itemize}
				\item Successful changes indicate control
				\item Unplanned downtime shows failures
				\item Meetings and budgets are inputs
			\end{itemize}
		\end{block}
		\begin{alertblock}{Critical Trap}
			Success rate, not volume, measures effectiveness.
		\end{alertblock}
	\end{frame}
	
	\begin{frame}{Q112: Change Freeze}
		\begin{block}{Question}
			What is the purpose of a change freeze?
		\end{block}
		\begin{enumerate}[(A)]
			\item To block all progress
			\item To reduce risk during critical business periods
			\item To save money
			\item To punish developers
		\end{enumerate}
	\end{frame}
	
	\begin{frame}{Q112: Explanation}
		\begin{block}{Correct Answer: B}
			\textbf{Freeze protects high-risk periods:}
			\begin{itemize}
				\item Financial close, holidays, peak seasons
				\item Reduces risk of disruption
				\item Emergency changes may still be approved
			\end{itemize}
		\end{block}
		\begin{alertblock}{Critical Trap}
			Freeze is a risk control, not a block on progress.
		\end{alertblock}
	\end{frame}
	
	\begin{frame}{Q113: Change vs Release}
		\begin{block}{Question}
			What is the MAIN distinction between change management and release management?
		\end{block}
		\begin{enumerate}[(A)]
			\item They are the same
			\item Change governs individual changes; release coordinates deployment packages
			\item Release is only for hardware
			\item Change is only for software
		\end{enumerate}
	\end{frame}
	
	\begin{frame}{Q113: Explanation}
		\begin{block}{Correct Answer: B}
			\textbf{Change governs; release coordinates:}
			\begin{itemize}
				\item Change evaluates and approves
				\item Release coordinates deployment
				\item Both reduce risk
			\end{itemize}
		\end{block}
		\begin{alertblock}{Critical Trap}
			Change approves; release deploys.
		\end{alertblock}
	\end{frame}
	
	\begin{frame}{Q114: Root Cause Analysis}
		\begin{block}{Question}
			Which tool BEST supports root cause analysis?
		\end{block}
		\begin{enumerate}[(A)]
			\item Gantt chart
			\item Ishikawa (fishbone) diagram
			\item Org chart
			\item Budget spreadsheet
		\end{enumerate}
	\end{frame}
	
	\begin{frame}{Q114: Explanation}
		\begin{block}{Correct Answer: B}
			\textbf{Fishbone finds causes:}
			\begin{itemize}
				\item Categorizes potential causes
				\item Structures problem-solving
				\item Supports problem management
			\end{itemize}
		\end{block}
		\begin{alertblock}{Critical Trap}
			Fishbone is for causes; Gantt is for scheduling.
		\end{alertblock}
	\end{frame}
	
	\begin{frame}{Q115: ITIL Service Value}
		\begin{block}{Question}
			What does ITIL focus on in IT operations?
		\end{block}
		\begin{enumerate}[(A)]
			\item Bypassing governance
			\item Delivering value through IT services aligned to business needs
			\item Only writing code
			\item Only managing hardware
		\end{enumerate}
	\end{frame}
	
	\begin{frame}{Q115: Explanation}
		\begin{block}{Correct Answer: B}
			\textbf{ITIL is about service value:}
			\begin{itemize}
				\item Services aligned to business
				\item Value realized through outcomes
				\item Continuous improvement
			\end{itemize}
		\end{block}
		\begin{alertblock}{Critical Trap}
			ITIL is about value, not just process.
		\end{alertblock}
	\end{frame}
	
	\begin{frame}{Q116: Incident Severity}
		\begin{block}{Question}
			Which factor should MOST drive incident severity classification?
		\end{block}
		\begin{enumerate}[(A)]
			\item Number of people who reported it
			\item Business impact and urgency
			\item Age of the ticket
			\item Vendor preference
		\end{enumerate}
	\end{frame}
	
	\begin{frame}{Q116: Explanation}
		\begin{block}{Correct Answer: B}
			\textbf{Impact and urgency drive severity:}
			\begin{itemize}
				\item Severity reflects business impact
				\item Urgency reflects time sensitivity
				\item Reports and age are secondary
			\end{itemize}
		\end{block}
		\begin{alertblock}{Critical Trap}
			Severity is impact plus urgency, not ticket count.
		\end{alertblock}
	\end{frame}
	
	\begin{frame}{Q117: Configuration Item}
		\begin{block}{Question}
			What is a configuration item in ITIL terms?
		\end{block}
		\begin{enumerate}[(A)]
			\item A budget line
			\item Any component under configuration management that is needed to deliver a service
			\item A user
			\item A vendor
		\end{enumerate}
	\end{frame}
	
	\begin{frame}{Q117: Explanation}
		\begin{block}{Correct Answer: B}
			\textbf{CIs are managed components:}
			\begin{itemize}
				\item Hardware, software, documentation
				\item Managed through the CMDB
				\item Enable impact analysis
			\end{itemize}
		\end{block}
		\begin{alertblock}{Critical Trap}
			CIs are managed components, not budgets.
		\end{alertblock}
	\end{frame}
	
	\begin{frame}{Q118: CMDB Purpose}
		\begin{block}{Question}
			What is the MAIN purpose of a Configuration Management Database?
		\end{block}
		\begin{enumerate}[(A)]
			\item Store employee records
			\item Maintain relationships between configuration items and support impact analysis
			\item Manage budgets
			\item Track marketing
		\end{enumerate}
	\end{frame}
	
	\begin{frame}{Q118: Explanation}
		\begin{block}{Correct Answer: B}
			\textbf{CMDB maps relationships:}
			\begin{itemize}
				\item Supports impact analysis
				\item Reveals dependencies
				\item Guides change risk
			\end{itemize}
		\end{block}
		\begin{alertblock}{Critical Trap}
			CMDB is about relationships, not records.
		\end{alertblock}
	\end{frame}
	
	\begin{frame}{Q119: Service Level Agreement}
		\begin{block}{Question}
			What is the primary purpose of a Service Level Agreement?
		\end{block}
		\begin{enumerate}[(A)]
			\item Define penalties only
			\item Define agreed service levels and expectations between provider and customer
			\item Replace contracts
			\item Avoid monitoring
		\end{enumerate}
	\end{frame}
	
	\begin{frame}{Q119: Explanation}
		\begin{block}{Correct Answer: B}
			\textbf{SLA defines expectations:}
			\begin{itemize}
				\item Service levels and targets
				\item Measurable expectations
				\item Supports governance
			\end{itemize}
		\end{block}
		\begin{alertblock}{Critical Trap}
			SLA is about expectations, not just penalties.
		\end{alertblock}
	\end{frame}
	
	\begin{frame}{Q120: Operational Metrics}
		\begin{block}{Question}
			Which metric BEST reflects operational stability?
		\end{block}
		\begin{enumerate}[(A)]
			\item Number of meetings
			\item Mean time to resolve incidents and system uptime
			\item Training slides
			\item Vendor discounts
		\end{enumerate}
	\end{frame}
	
	\begin{frame}{Q120: Explanation}
		\begin{block}{Correct Answer: B}
			\textbf{Stability is MTTR and uptime:}
			\begin{itemize}
				\item MTTR shows response effectiveness
				\item Uptime shows availability
				\item Both reflect operations
			\end{itemize}
		\end{block}
		\begin{alertblock}{Critical Trap}
			Stability is uptime and resolution, not process volume.
		\end{alertblock}
	\end{frame}
	
	% ============================================================
	\section{3. Enterprise Resiliency}
	% ============================================================
	
	\begin{frame}{Q121: BIA vs DR}
		\begin{block}{Question}
			What is the role of a Business Impact Analysis in disaster recovery?
		\end{block}
		\begin{enumerate}[(A)]
			\item Unrelated
			\item BIA defines critical processes and recovery objectives that DR implements
			\item DR replaces BIA
			\item BIA replaces BC
		\end{enumerate}
	\end{frame}
	
	\begin{frame}{Q121: Explanation}
		\begin{block}{Correct Answer: B}
			\textbf{BIA informs DR:}
			\begin{itemize}
				\item BIA defines priorities and RTO/RPO
				\item DR implements recovery
				\item BIA is the input, DR is the execution
			\end{itemize}
		\end{block}
		\begin{alertblock}{Critical Trap}
			BIA comes before DR; DR implements BIA.
		\end{alertblock}
	\end{frame}
	
	\begin{frame}{Q122: RTO Realism}
		\begin{block}{Question}
			An RTO of 5 minutes is proposed for a system with hourly backups. What is the issue?
		\end{block}
		\begin{enumerate}[(A)]
			\item No issue
			\item RPO does not support RTO
			\item RTO is too long
			\item Backups are too frequent
		\end{enumerate}
	\end{frame}
	
	\begin{frame}{Q122: Explanation}
		\begin{block}{Correct Answer: B}
			\textbf{RPO must support RTO:}
			\begin{itemize}
				\item Hourly backups cannot restore quickly
				\item RTO must match recovery technology
				\item Realistic objectives matter
			\end{itemize}
		\end{block}
		\begin{alertblock}{Critical Trap}
			RTO must be realistic against backup and recovery capability.
		\end{alertblock}
	\end{frame}
	
	\begin{frame}{Q123: Cold vs Warm vs Hot Site}
		\begin{block}{Question}
			Which site type offers the fastest recovery?
		\end{block}
		\begin{enumerate}[(A)]
			\item Cold site
			\item Warm site
			\item Hot site
			\item No site
		\end{enumerate}
	\end{frame}
	
	\begin{frame}{Q123: Explanation}
		\begin{block}{Correct Answer: C}
			\textbf{Hot site is fastest:}
			\begin{itemize}
				\item Fully operational duplicate
				\item Immediate failover
				\item Highest cost
			\end{itemize}
		\end{block}
		\begin{alertblock}{Critical Trap}
			Hot site is fastest, costliest.
		\end{alertblock}
	\end{frame}
	
	\begin{frame}{Q124: Cold Site Use}
		\begin{block}{Question}
			When is a cold site MOST appropriate?
		\end{block}
		\begin{enumerate}[(A)]
			\item Mission-critical 24/7 systems
			\item Low-criticality systems with long tolerable downtime
			\item Trading systems
			\item Always
		\end{enumerate}
	\end{frame}
	
	\begin{frame}{Q124: Explanation}
		\begin{block}{Correct Answer: B}
			\textbf{Cold sites suit long RTOs:}
			\begin{itemize}
				\item No hardware pre-installed
				\item Requires time to build
				\item Suitable for low criticality
			\end{itemize}
		\end{block}
		\begin{alertblock}{Critical Trap}
			Cold sites have long recovery times.
		\end{alertblock}
	\end{frame}
	
	\begin{frame}{Q125: Backup vs Replication}
		\begin{block}{Question}
			What is the MAIN difference between backup and replication?
		\end{block}
		\begin{enumerate}[(A)]
			\item They are identical
			\item Backup copies data periodically; replication mirrors data continuously or near-continuously
			\item Backup is continuous; replication is periodic
			\item Both are manual
		\end{enumerate}
	\end{frame}
	
	\begin{frame}{Q125: Explanation}
		\begin{block}{Correct Answer: B}
			\textbf{Backup periodic; replication continuous:}
			\begin{itemize}
				\item Backup creates recovery points
				\item Replication maintains mirrors
				\item Both reduce recovery time
			\end{itemize}
		\end{block}
		\begin{alertblock}{Critical Trap}
			Backup is periodic; replication mirrors.
		\end{alertblock}
	\end{frame}
	
	\begin{frame}{Q126: Ransomware Recovery}
		\begin{block}{Question}
			Which is MOST important for ransomware recovery?
		\end{block}
		\begin{enumerate}[(A)]
			\item Number of antivirus tools
			\item Immutable, offline backups and tested restore capability
			\item Faster email
			\item More vendors
		\end{enumerate}
	\end{frame}
	
	\begin{frame}{Q126: Explanation}
		\begin{block}{Correct Answer: B}
			\textbf{Immutable backups protect against ransomware:}
			\begin{itemize}
				\item Online backups can be encrypted
				\item Immutable and offline backups stay clean
				\item Restore testing is critical
			\end{itemize}
		\end{block}
		\begin{alertblock}{Critical Trap}
			Ransomware encrypts connected backups. Immutable backups matter.
		\end{alertblock}
	\end{frame}
	
	\begin{frame}{Q127: Recovery Prioritization}
		\begin{block}{Question}
			Which system should be restored FIRST after a disaster?
		\end{block}
		\begin{enumerate}[(A)]
			\item Lowest-cost system
			\item System with the smallest RTO critical to business continuity
			\item Oldest system
			\item Newest system
		\end{enumerate}
	\end{frame}
	
	\begin{frame}{Q127: Explanation}
		\begin{block}{Correct Answer: B}
			\textbf{Priority follows RTO and criticality:}
			\begin{itemize}
				\item Shortest RTO first
				\item Critical business processes
				\item Recovery sequence matters
			\end{itemize}
		\end{block}
		\begin{alertblock}{Critical Trap}
			Recovery priority = criticality and RTO, not cost or age.
		\end{alertblock}
	\end{frame}
	
	\begin{frame}{Q128: BCP Communication}
		\begin{block}{Question}
			Which is essential for BCP communication during a disruption?
		\end{block}
		\begin{enumerate}[(A)]
			\item Marketing brochures
			\item Predefined channels, roles, and templates for internal and external stakeholders
			\item Only internal chat
			\item Only email
		\end{enumerate}
	\end{frame}
	
	\begin{frame}{Q128: Explanation}
		\begin{block}{Correct Answer: B}
			\textbf{Preplanned communication saves time:}
			\begin{itemize}
				\item Defined channels
				\item Named roles
				\item Ready templates
			\end{itemize}
		\end{block}
		\begin{alertblock}{Critical Trap}
			Communication must be planned, not improvised.
		\end{alertblock}
	\end{frame}
	
	\begin{frame}{Q129: BCP Exercise Types}
		\begin{block}{Question}
			Which exercise type provides the HIGHEST confidence at the LOWEST cost?
		\end{block}
		\begin{enumerate}[(A)]
			\item Full failover with real services
			\item Tabletop exercise with key stakeholders
			\item Ignoring the plan
			\item Adding more tools
		\end{enumerate}
	\end{frame}
	
	\begin{frame}{Q129: Explanation}
		\begin{block}{Correct Answer: B}
			\textbf{Tabletop balances cost and value:}
			\begin{itemize}
				\item Low cost
				\item Reveals gaps in roles and communications
				\item Prepares for full tests
			\end{itemize}
		\end{block}
		\begin{alertblock}{Critical Trap}
			Tabletop = highest value at low cost.
		\end{alertblock}
	\end{frame}
	
	\begin{frame}{Q130: BCP Update Triggers}
		\begin{block}{Question}
			Which event should trigger a BCP update?
		\end{block}
		\begin{enumerate}[(A)]
			\item Staff birthday
			\item Major business or technology change
			\item Vendor price change
			\item Office decoration
		\end{enumerate}
	\end{frame}
	
	\begin{frame}{Q130: Explanation}
		\begin{block}{Correct Answer: B}
			\textbf{Significant change triggers review:}
			\begin{itemize}
				\item Architecture or process changes
				\item Affects RTO/RPO
				\item Keeps BCP current
			\end{itemize}
		\end{block}
		\begin{alertblock}{Critical Trap}
			BCP updates follow significant change.
		\end{alertblock}
	\end{frame}
	
	% ============================================================
	\section{4. Data Lifecycle Management}
	% ============================================================
	
	\begin{frame}{Q131: Data Classification Levels}
		\begin{block}{Question}
			Which is a typical data classification set?
		\end{block}
		\begin{enumerate}[(A)]
			\item Fast, slow, medium
			\item Public, internal, confidential, restricted
			\item Red, green, blue
			\item Old, new, older
		\end{enumerate}
	\end{frame}
	
	\begin{frame}{Q131: Explanation}
		\begin{block}{Correct Answer: B}
			\textbf{Public, internal, confidential, restricted:}
			\begin{itemize}
				\item Common classification levels
				\item Sensitivity determines controls
				\item Aligns with regulations
			\end{itemize}
		\end{block}
		\begin{alertblock}{Critical Trap}
			Classification is sensitivity-based, not time-based.
		\end{alertblock}
	\end{frame}
	
	\begin{frame}{Q132: Data Ownership vs Custody}
		\begin{block}{Question}
			Who decides who can access specific data?
		\end{block}
		\begin{enumerate}[(A)]
			\item The custodian
			\item The data owner
			\item The user
			\item The vendor
		\end{enumerate}
	\end{frame}
	
	\begin{frame}{Q132: Explanation}
		\begin{block}{Correct Answer: B}
			\textbf{Data owners authorize access:}
			\begin{itemize}
				\item Owners determine access rights
				\item Custodians implement controls
				\item Users consume data
			\end{itemize}
		\end{block}
		\begin{alertblock}{Critical Trap}
			Owners authorize; custodians implement.
		\end{alertblock}
	\end{frame}
	
	\begin{frame}{Q133: Retention vs Legal Hold}
		\begin{block}{Question}
			A retention schedule says destroy after 7 years. A legal hold is in place. What applies?
		\end{block}
		\begin{enumerate}[(A)]
			\item Destroy per schedule
			\item Legal hold overrides the retention schedule
			\item Ignore both
			\item Destroy immediately
		\end{enumerate}
	\end{frame}
	
	\begin{frame}{Q133: Explanation}
		\begin{block}{Correct Answer: B}
			\textbf{Legal hold overrides:}
			\begin{itemize}
				\item Preserve during hold
				\item Destruction during hold is a violation
				\item Release the hold when cleared
			\end{itemize}
		\end{block}
		\begin{alertblock}{Critical Trap}
			Legal holds suspend retention rules.
		\end{alertblock}
	\end{frame}
	
	\begin{frame}{Q134: DLP Enforcement}
		\begin{block}{Question}
			Which data classification MOST benefits from DLP enforcement?
		\end{block}
		\begin{enumerate}[(A)]
			\item Public data
			\item Highly confidential data
			\item Marketing brochures
			\item Training materials
		\end{enumerate}
	\end{frame}
	
	\begin{frame}{Q134: Explanation}
		\begin{block}{Correct Answer: B}
			\textbf{DLP protects sensitive data:}
			\begin{itemize}
				\item Confidential and restricted classes
				\item DLP prevents leakage
				\item Public data has less need
			\end{itemize}
		\end{block}
		\begin{alertblock}{Critical Trap}
			DLP is for high-sensitivity data.
		\end{alertblock}
	\end{frame}
	
	\begin{frame}{Q135: Data Lifecycle Creation}
		\begin{block}{Question}
			At which lifecycle stage should classification occur?
		\end{block}
		\begin{enumerate}[(A)]
			\item Creation
			\item Archiving
			\item Destruction
			\item Only after a year
		\end{enumerate}
	\end{frame}
	
	\begin{frame}{Q135: Explanation}
		\begin{block}{Correct Answer: A}
			\textbf{Classify at creation:}
			\begin{itemize}
				\item Classify early to enforce controls
				\item Late classification is often forgotten
				\item Creation is the start of protection
			\end{itemize}
		\end{block}
		\begin{alertblock}{Critical Trap}
			Classify early, protect early.
		\end{alertblock}
	\end{frame}
	
	\begin{frame}{Q136: Data Minimization GDPR}
		\begin{block}{Question}
			Which principle aligns with GDPR data minimization?
		\end{block}
		\begin{enumerate}[(A)]
			\item Collect everything
			\item Collect only data needed for stated purposes
			\item Retain data forever
			\item Ignore privacy
		\end{enumerate}
	\end{frame}
	
	\begin{frame}{Q136: Explanation}
		\begin{block}{Correct Answer: B}
			\textbf{Collect only what is needed:}
			\begin{itemize}
				\item Reduces exposure
				\item Supports privacy
				\item Improves compliance
			\end{itemize}
		\end{block}
		\begin{alertblock}{Critical Trap}
			Data minimization = need-to-know collection.
		\end{alertblock}
	\end{frame}
	
	\begin{frame}{Q137: Data Destruction Certificate}
		\begin{block}{Question}
			Why is a Certificate of Destruction important?
		\end{block}
		\begin{enumerate}[(A)]
			\item Marketing
			\item Provides documented evidence that data was securely destroyed
			\item Tax filing
			\item Vendor discounts
		\end{enumerate}
	\end{frame}
	
	\begin{frame}{Q137: Explanation}
		\begin{block}{Correct Answer: B}
			\textbf{Evidence matters:}
			\begin{itemize}
				\item Documents compliant destruction
				\item Supports audit
				\item Protects legally
			\end{itemize}
		\end{block}
		\begin{alertblock}{Critical Trap}
			Destroyed without evidence is not defensible.
		\end{alertblock}
	\end{frame}
	
	\begin{frame}{Q138: Cloud Data Location}
		\begin{block}{Question}
			Why is data location important for cloud-stored data?
		\end{block}
		\begin{enumerate}[(A)]
			\item Latency only
			\item Legal jurisdiction and data residency requirements
			\item Vendor pricing
			\item Logo placement
		\end{enumerate}
	\end{frame}
	
	\begin{frame}{Q138: Explanation}
		\begin{block}{Correct Answer: B}
			\textbf{Location determines legal exposure:}
			\begin{itemize}
				\item Data residency laws
				\item Jurisdiction impacts regulator access
				\item Compliance risks
			\end{itemize}
		\end{block}
		\begin{alertblock}{Critical Trap}
			Data location is a legal concern.
		\end{alertblock}
	\end{frame}
	
	\begin{frame}{Q139: Data Subject Request}
		\begin{block}{Question}
			What should an organization do FIRST when receiving a data subject access request?
		\end{block}
		\begin{enumerate}[(A)]
			\item Delete data
			\item Verify the identity of the requester
			\item Forward to marketing
			\item Ignore it
		\end{enumerate}
	\end{frame}
	
	\begin{frame}{Q139: Explanation}
		\begin{block}{Correct Answer: B}
			\textbf{Verify identity first:}
			\begin{itemize}
				\item Prevents unauthorized disclosure
				\item Ensures right person
				\item Then fulfill the request
			\end{itemize}
		\end{block}
		\begin{alertblock}{Critical Trap}
			Verify identity before responding to any DSR.
		\end{alertblock}
	\end{frame}
	
	\begin{frame}{Q140: Archival Encryption}
		\begin{block}{Question}
			Why should archived data remain encrypted?
		\end{block}
		\begin{enumerate}[(A)]
			\item To save cost
			\item To protect against exposure if storage is compromised
			\item To speed up retrieval
			\item To avoid retention
		\end{enumerate}
	\end{frame}
	
	\begin{frame}{Q140: Explanation}
		\begin{block}{Correct Answer: B}
			\textbf{Archived data remains at risk:}
			\begin{itemize}
				\item Long-term storage
				\item Sensitive data persists
				\item Encryption reduces exposure
			\end{itemize}
		\end{block}
		\begin{alertblock}{Critical Trap}
			Archived data still needs protection.
		\end{alertblock}
	\end{frame}
	
	% ============================================================
	\section{5. SDLC and Emerging Technologies}
	% ============================================================
	
	\begin{frame}{Q141: Waterfall vs Agile Risk}
		\begin{block}{Question}
			Which model has GREATER change-related risk?
		\end{block}
		\begin{enumerate}[(A)]
			\item Waterfall
			\item Agile
			\item Both equal
			\item Neither}
		{B}
		\begin{itemize}
			\item
			\item Agile changes rapidly
			\item Continuous testing is essential
			\item Waterfall changes are costly
		\end{itemize}
		\begin{alertblock}{Critical Trap}
			Agile risk = change management.
		\end{alertblock}
	\end{frame}
	
	\begin{frame}{Q142: Secure Code Review}
		\begin{block}{Question}
			What is the PRIMARY benefit of secure code review?
		\end{block}
		\begin{enumerate}[(A)]
			\item Faster delivery
			\item Early detection of vulnerabilities before release
			\item Lower storage
			\item More features
		\end{enumerate}
	\end{frame}
	
	\begin{frame}{Q142: Explanation}
		\begin{block}{Correct Answer: B}
			\textbf{Early detection reduces cost:}
			\begin{itemize}
				\item Defects cheaper to fix early
				\item Security reviewed before exposure
				\item Complements testing
			\end{itemize}
		\end{block}
		\begin{alertblock}{Critical Trap}
			Early detection is cheaper than late repair.
		\end{alertblock}
	\end{frame}
	
	\begin{frame}{Q143: Dependency Scanning}
		\begin{block}{Question}
			Why is dependency scanning critical in modern SDLC?
		\end{block}
		\begin{enumerate}[(A)]
			\item Vendors require it
			\item Third-party libraries introduce vulnerabilities
			\item It speeds deployment
			\item It reduces costs
		\end{enumerate}
	\end{frame}
	
	\begin{frame}{Q143: Explanation}
		\begin{block}{Correct Answer: B}
			\textbf{Dependencies introduce risk:}
			\begin{itemize}
				\item Open-source libraries can have CVEs
				\item Regular scanning detects them
				\item Reduces supply chain risk
			\end{itemize}
		\end{block}
		\begin{alertblock}{Critical Trap}
			Dependency risk is a growing threat vector.
		\end{alertblock}
	\end{frame}
	
	\begin{frame}{Q144: Container Security}
		\begin{block}{Question}
			Which security control is MOST important for containerized workloads?
		\end{block}
		\begin{enumerate}[(A)]
			\item Antivirus on hosts
			\item Image scanning and runtime monitoring
			\item Only firewalls
			\item Only encryption
		\end{enumerate}
	\end{frame}
	
	\begin{frame}{Q144: Explanation}
		\begin{block}{Correct Answer: B}
			\textbf{Image scanning plus runtime:}
			\begin{itemize}
				\item Scans for vulnerable images
				\item Runtime detection
				\item Covers container lifecycle
			\end{itemize}
		\end{block}
		\begin{alertblock}{Critical Trap}
			Containers need image and runtime controls.
		\end{alertblock}
	\end{frame}
	
	\begin{frame}{Q145: Cloud Shared Responsibility}
		\begin{block}{Question}
			What does the shared responsibility model mean in cloud?
		\end{block}
		\begin{enumerate}[(A)]
			\item Provider responsible for everything
			\item Provider secures infrastructure; customer secures configuration, data, and identity
			\item Customer responsible for everything
			\item No responsibility
		\end{enumerate}
	\end{frame}
	
	\begin{frame}{Q145: Explanation}
		\begin{block}{Correct Answer: B}
			\textbf{Security is shared:}
			\begin{itemize}
				\item Provider secures physical and hypervisor
				\item Customer secures data and configuration
				\item Misunderstanding causes breaches
			\end{itemize}
		\end{block}
		\begin{alertblock}{Critical Trap}
			Cloud security is shared, not delegated.
		\end{alertblock}
	\end{frame}
	
	\begin{frame}{Q146: SaaS Tenant Isolation}
		\begin{block}{Question}
			Why is SaaS tenant isolation important?
		\end{block}
		\begin{enumerate}[(A)]
			\item Performance only
			\item Prevents data leakage between customers
			\item Reduces cost
			\item Simplifies audits
		\end{enumerate}
	\end{frame}
	
	\begin{frame}{Q146: Explanation}
		\begin{block}{Correct Answer: B}
			\textbf{Isolation prevents cross-tenant leakage:}
			\begin{itemize}
				\item Multi-tenancy risk
				\item Isolation required
				\item Verify in contracts
			\end{itemize}
		\end{block}
		\begin{alertblock}{Critical Trap}
			Tenant isolation is a critical SaaS control.
		\end{alertblock}
	\end{frame}
	
	\begin{frame}{Q147: Quantum Safe Cryptography}
		\begin{block}{Question}
			What should organizations do about post-quantum cryptography?
		\end{block}
		\begin{enumerate}[(A)]
			\item Wait until quantum computers exist
			\item Plan migration to quantum-safe algorithms
			\item Ignore
			\item Use only symmetric
		\end{enumerate}
	\end{frame}
	
	\begin{frame}{Q147: Explanation}
		\begin{block}{Correct Answer: B}
			\textbf{Plan migration now:}
			\begin{itemize}
				\item Harvest-now-decrypt-later risk
				\item Post-quantum algorithms available
				\item Migration takes years
			\end{itemize}
		\end{block}
		\begin{alertblock}{Critical Trap}
			Quantum migration must start early.
		\end{alertblock}
	\end{frame}
	
	\begin{frame}{Q148: IoT Segmentation}
		\begin{block}{Question}
			Why segment IoT devices from corporate networks?
		\end{block}
		\begin{enumerate}[(A)]
			\item Better performance
			\item Limits blast radius if devices are compromised
			\item Reduces cost
			\item Simpler backups
		\end{enumerate}
	\end{frame}
	
	\begin{frame}{Q148: Explanation}
		\begin{block}{Correct Answer: B}
			\textbf{Segmentation limits damage:}
			\begin{itemize}
				\item IoT often has weak security
				\item Segmentation contains compromise
				\item Reduces lateral movement
			\end{itemize}
		\end{block}
		\begin{alertblock}{Critical Trap}
			IoT segmentation is a mitigation, not a fix.
		\end{alertblock}
	\end{frame}
	
	\begin{frame}{Q149: AI Explainability}
		\begin{block}{Question}
			Why is AI explainability important for governance?
		\end{block}
		\begin{enumerate}[(A)]
			\item Faster training
			\item Enables oversight, audit, and accountability
			\item Cheaper GPUs
			\item Marketing value
		\end{enumerate}
	\end{frame}
	
	\begin{frame}{Q149: Explanation}
		\begin{block}{Correct Answer: B}
			\textbf{Explainability enables governance:}
			\begin{itemize}
				\item Decisions can be audited
				\item Bias can be detected
				\item Accountability supported
			\end{itemize}
		\end{block}
		\begin{alertblock}{Critical Trap}
			Opaque AI cannot be governed.
		\end{alertblock}
	\end{frame}
	
	\begin{frame}{Q150: Smart Contract Risk}
		\begin{block}{Question}
			What is a key governance risk with smart contracts?
		\end{block}
		\begin{enumerate}[(A)]
			\item Slow execution
			\item Bugs cannot be easily corrected once deployed
			\item Too many users
			\item High cost
		\end{enumerate}
	\end{frame}
	
	\begin{frame}{Q150: Explanation}
		\begin{block}{Correct Answer: B}
			\textbf{Immutability has consequences:}
			\begin{itemize}
				\item Deployed contracts immutable
				\item Bugs persist
				\item Requires thorough testing before deployment
			\end{itemize}
		\end{block}
		\begin{alertblock}{Critical Trap}
			Smart contract bugs are permanent.
		\end{alertblock}
	\end{frame}
	
	% ============================================================
	\section{6. Information Security Principles}
	% ============================================================
	
	\begin{frame}{Q151: IAAA Accountability}
		\begin{block}{Question}
			Which principle logs and monitors what users did?
		\end{block}
		\begin{enumerate}[(A)]
			\item Identification
			\item Authentication
			\item Authorization
			\item Accountability
		\end{enumerate}
	\end{frame}
	
	\begin{frame}{Q151: Explanation}
		\begin{block}{Correct Answer: D}
			\textbf{Accountability tracks actions:}
			\begin{itemize}
				\item Logs record activity
				\item Supports investigation
				\item Non-repudiation support
			\end{itemize}
		\end{block}
		\begin{alertblock}{Critical Trap}
			Accountability = logging and monitoring.
		\end{alertblock}
	\end{frame}
	
	\begin{frame}{Q152: Authorization Definition}
		\begin{block}{Question}
			Authorization defines:
		\end{block}
		\begin{enumerate}[(A)]
			\item Who you are
			\item What you can do
			\item How you authenticate
			\item Where you log in
		\end{enumerate}
	\end{frame}
	
	\begin{frame}{Q152: Explanation}
		\begin{block}{Correct Answer: B}
			\textbf{Authorization is permission:}
			\begin{itemize}
				\item Determines actions allowed
				\item Enforced through access control
				\item Post-authentication decision
			\end{itemize}
		\end{block}
		\begin{alertblock}{Critical Trap}
			Authorization = what you can do.
		\end{alertblock}
	\end{frame}
	
	\begin{frame}{Q153: Hashing Password Storage}
		\begin{block}{Question}
			Why store passwords as hashes instead of encryption?
		\end{block}
		\begin{enumerate}[(A)]
			\item Faster
			\item Hashes cannot be reversed to plaintext
			\item Cheaper
			\item Shorter
		\end{enumerate}
	\end{frame}
	
	\begin{frame}{Q153: Explanation}
		\begin{block}{Correct Answer: B}
			\textbf{Hashes are one-way:}
			\begin{itemize}
				\item Reversal infeasible
				\item Even if breached, plaintext is not exposed
				\item Salt adds further protection
			\end{itemize}
		\end{block}
		\begin{alertblock}{Critical Trap}
			Hashing is one-way; encryption is reversible.
		\end{alertblock}
	\end{frame}
	
	\begin{frame}{Q154: Salt Purpose}
		\begin{block}{Question}
			What is the purpose of salting password hashes?
		\end{block}
		\begin{enumerate}[(A)]
			\item Speeds up login
			\item Prevents identical passwords from producing the same hash
			\item Reduces storage
			\item Simplifies resets
		\end{enumerate}
	\end{frame}
	
	\begin{frame}{Q154: Explanation}
		\begin{block}{Correct Answer: B}
			\textbf{Salt makes hashes unique:}
			\begin{itemize}
				\item Different salt per password
				\item Defeats rainbow tables
				\item Strengthens storage
			\end{itemize}
		\end{block}
		\begin{alertblock}{Critical Trap}
			Salting defeats precomputed attacks.
		\end{alertblock}
	\end{frame}
	
	\begin{frame}{Q155: TLS Purpose}
		\begin{block}{Question}
			What does TLS provide in transit?
		\end{block}
		\begin{enumerate}[(A)]
			\item Storage encryption only
			\item Confidentiality and integrity over the network
			\item Antivirus
			\item Backups
		\end{enumerate}
	\end{frame}
	
	\begin{frame}{Q155: Explanation}
		\begin{block}{Correct Answer: B}
			\textbf{TLS protects transit:}
			\begin{itemize}
				\item Encrypts traffic
				\item Verifies integrity
				\item Certificate-based trust
			\end{itemize}
		\end{block}
		\begin{alertblock}{Critical Trap}
			TLS is transit protection, not storage.
		\end{alertblock}
	\end{frame}
	
	\begin{frame}{Q156: Firewall Purpose}
		\begin{block}{Question}
			What is the PRIMARY purpose of a firewall?
		\end{block}
		\begin{enumerate}[(A)]
			\item Encrypt traffic
			\item Enforce network policy on traffic
			\item Back up data
			\item Detect malware
		\end{enumerate}
	\end{frame}
	
	\begin{frame}{Q156: Explanation}
		\begin{block}{Correct Answer: B}
			\textbf{Firewalls enforce policy:}
			\begin{itemize}
				\item Allow or deny traffic
				\item Rule-based
				\item Complements other controls
			\end{itemize}
		\end{block}
		\begin{alertblock}{Critical Trap}
			Firewalls enforce policy; they do not encrypt.
		\end{alertblock}
	\end{frame}
	
	\begin{frame}{Q157: IDS vs IPS}
		\begin{block}{Question}
			Which is passive and which is active?
		\end{block}
		\begin{enumerate}[(A)]
			\item IDS passive; IPS active
			\item IDS active; IPS passive
			\item Both active
			\item Both passive}
		{A}
		\begin{itemize}
			\item
			\item IDS detects only
			\item IPS detects and blocks
			\item Choose based on requirements
		\end{itemize}
		\begin{alertblock}{Critical Trap}
			IDS detects; IPS blocks.
		\end{alertblock}
	\end{frame}
	
	\begin{frame}{Q158: WAF Purpose}
		\begin{block}{Question}
			What does a Web Application Firewall primarily do?
		\end{block}
		\begin{enumerate}[(A)]
			\item Encrypt web content
			\item Filter HTTP traffic to protect web apps from common attacks
			\item Back up websites
			\item Manage DNS
		\end{enumerate}
	\end{frame}
	
	\begin{frame}{Q158: Explanation}
		\begin{block}{Correct Answer: B}
			\textbf{WAF filters HTTP:}
			\begin{itemize}
				\item Protects against OWASP Top 10
				\item Complements firewalls
				\item Layer 7 filtering
			\end{itemize}
		\end{block}
		\begin{alertblock}{Critical Trap}
			WAF is layer 7 web protection.
		\end{alertblock}
	\end{frame}
	
	\begin{frame}{Q159: Zero Trust Identity}
		\begin{block}{Question}
			What identity principle is central to Zero Trust?
		\end{block}
		\begin{enumerate}[(A)]
			\item Trust all internal users
			\item Verify every access request regardless of location
			\item Use only passwords
			\item Trust all vendors
		\end{enumerate}
	\end{frame}
	
	\begin{frame}{Q159: Explanation}
		\begin{block}{Correct Answer: B}
			\textbf{Verify every request:}
			\begin{itemize}
				\item Identity and context
				\item Continuous verification
				\item No implicit trust
			\end{itemize}
		\end{block}
		\begin{alertblock}{Critical Trap}
			Zero Trust = verify every access, every time.
		\end{alertblock}
	\end{frame}
	
	\begin{frame}{Q160: MFA Fatigue Attack}
		\begin{block}{Question}
			What is MFA fatigue and how is it mitigated?
		\end{block}
		\begin{enumerate}[(A)]
			\item Users forget passwords; reset them
			\item Attackers spam MFA prompts; mitigate with number matching and rate limits
			\item MFA never fails
			\item MFA is unused
		\end{enumerate}
	\end{frame}
	
	\begin{frame}{Q160: Explanation}
		\begin{block}{Correct Answer: B}
			\textbf{Spamming prompts tricks users:}
			\begin{itemize}
				\item Attackers flood MFA prompts
				\item Users approve by accident
				\item Number matching mitigates
			\end{itemize}
		\end{block}
		\begin{alertblock}{Critical Trap}
			MFA fatigue attacks exploit user fatigue.
		\end{alertblock}
	\end{frame}
	
	\begin{frame}{Q161: Identity Provider Compromise}
		\begin{block}{Question}
			What is the PRIMARY risk if the identity provider is compromised?
		\end{block}
		\begin{enumerate}[(A)]
			\item Slower logins
			\item Attacker can impersonate users and access resources
			\item Loss of backups
			\item Slower printing
		\end{enumerate}
	\end{frame}
	
	\begin{frame}{Q161: Explanation}
		\begin{block}{Correct Answer: B}
			\textbf{IdP compromise is catastrophic:}
			\begin{itemize}
				\item Central trust point
				\item Attackers can impersonate users
				\item Protect with MFA, monitoring, least privilege
			\end{itemize}
		\end{block}
		\begin{alertblock}{Critical Trap}
			IdP is a high-value target.
		\end{alertblock}
	\end{frame}
	
	\begin{frame}{Q162: Cloud IAM Best Practice}
		\begin{block}{Question}
			What is a best practice for cloud IAM?
		\end{block}
		\begin{enumerate}[(A)]
			\item Use root account daily
			\item Use least privilege, MFA, and role-based access
			\item Share admin credentials
			\item Disable logging
		\end{enumerate}
	\end{frame}
	
	\begin{frame}{Q162: Explanation}
		\begin{block}{Correct Answer: B}
			\textbf{Least privilege, MFA, RBAC:}
			\begin{itemize}
				\item Root account should not be used daily
				\item Least privilege limits damage
				\item Logging is essential
			\end{itemize}
		\end{block}
		\begin{alertblock}{Critical Trap}
			Cloud IAM requires discipline, not convenience.
		\end{alertblock}
	\end{frame}
	
	\begin{frame}{Q163: Data Residency Risk}
		\begin{block}{Question}
			What risk does storing data in a foreign cloud region introduce?
		\end{block}
		\begin{enumerate}[(A)]
			\item Only latency
			\item Legal jurisdiction and regulatory requirements
			\item Slower backups
			\item Higher cost only
		\end{enumerate}
	\end{frame}
	
	\begin{frame}{Q163: Explanation}
		\begin{block}{Correct Answer: B}
			\textbf{Legal exposure follows data:}
			\begin{itemize}
				\item Different jurisdictions
				\item Data residency rules
				\item Regulator access varies
			\end{itemize}
		\end{block}
		\begin{alertblock}{Critical Trap}
			Data location brings legal exposure.
		\end{alertblock}
	\end{frame}
	
	\begin{frame}{Q164: Confidential Computing}
		\begin{block}{Question}
			What does confidential computing protect?
		\end{block}
		\begin{enumerate}[(A)]
			\item Only network traffic
			\item Data in use inside trusted execution environments
			\item Only backups
			\item Only usernames
		\end{enumerate}
	\end{frame}
	
	\begin{frame}{Q164: Explanation}
		\begin{block}{Correct Answer: B}
			\textbf{Protects data in use:}
			\begin{itemize}
				\item Trusted execution environments
				\item Protects from host-level exposure
				\item Complements encryption at rest and in transit
			\end{itemize}
		\end{block}
		\begin{alertblock}{Critical Trap}
			Confidential computing = data in use.
		\end{alertblock}
	\end{frame}
	
	\begin{frame}{Q165: Privileged Access Management}
		\begin{block}{Question}
			What is the main goal of Privileged Access Management?
		\end{block}
		\begin{enumerate}[(A)]
			\item Give everyone admin access
			\item Control, monitor, and audit use of privileged accounts
			\item Remove all admins
			\item Avoid logging
		\end{enumerate}
	\end{frame}
	
	\begin{frame}{Q165: Explanation}
		\begin{block}{Correct Answer: B}
			\textbf{PAM controls admin access:}
			\begin{itemize}
				\item Vault credentials
				\item Session monitoring
				\item Approval workflows
			\end{itemize}
		\end{block}
		\begin{alertblock}{Critical Trap}
			PAM protects the keys to the kingdom.
		\end{alertblock}
	\end{frame}
	
	\begin{frame}{Q166: Security Awareness Metrics}
		\begin{block}{Question}
			Which metric BEST shows security awareness effectiveness?
		\end{block}
		\begin{enumerate}[(A)]
			\item Completion rate
			\item Reduction in successful phishing clicks and increase in reported phishing
			\item Number of slides
			\item Training hours}
		{B}
		\begin{itemize}
			\item
			\item Behavior change matters
			\item Completion is activity
			\item Reports and click rates reflect behavior
		\end{itemize}
		\begin{alertblock}{Critical Trap}
			Behavior metrics matter more than completion.
		\end{alertblock}
	\end{frame}
	
	\begin{frame}{Q167: Insider Threat Indicators}
		\begin{block}{Question}
			Which is a red flag for insider threat?
		\end{block}
		\begin{enumerate}[(A)]
			\item Timely access requests
			\item Unusual access patterns, mass downloads, or access near resignation
			\item Regular training completion
			\item Using MFA
		\end{enumerate}
	\end{frame}
	
	\begin{frame}{Q167: Explanation}
		\begin{block}{Correct Answer: B}
			\textbf{Unusual patterns are red flags:}
			\begin{itemize}
				\item Mass data access
				\item Odd hours
				\item Access after notice
			\end{itemize}
		\end{block}
		\begin{alertblock}{Critical Trap}
			Insider threats show behavioral anomalies.
		\end{alertblock}
	\end{frame}
	
	\begin{frame}{Q168: Third-Party Breach Risk}
		\begin{block}{Question}
			When a third party suffers a breach of your data, who is regulated?
		\end{block}
		\begin{enumerate}[(A)]
			\item Only the third party
			\item Your organization as data controller
			\item Only the regulator
			\item No one
		\end{enumerate}
	\end{frame}
	
	\begin{frame}{Q168: Explanation}
		\begin{block}{Correct Answer: B}
			\textbf{Controllers remain regulated:}
			\begin{itemize}
				\item You are responsible to regulators
				\item Third parties act on your behalf
				\item Contracts and oversight matter
			\end{itemize}
		\end{block}
		\begin{alertblock}{Critical Trap}
			Outsourcing does not transfer regulatory accountability.
		\end{alertblock}
	\end{frame}
	
	\begin{frame}{Q169: Data Breach Root Cause}
		\begin{block}{Question}
			Which is a common root cause of data breaches?
		\end{block}
		\begin{enumerate}[(A)]
			\item Strong MFA everywhere
			\item Misconfiguration and weak access management
			\item Frequent training
			\item Regular audits
		\end{enumerate}
	\end{frame}
	
	\begin{frame}{Q169: Explanation}
		\begin{block}{Correct Answer: B}
			\textbf{Misconfigurations and access weaknesses:}
			\begin{itemize}
				\item Common breach causes
				\item Configuration management matters
				\item Access reviews reduce risk
			\end{itemize}
		\end{block}
		\begin{alertblock}{Critical Trap}
			Misconfiguration is a top breach cause.
		\end{alertblock}
	\end{frame}
	
	\begin{frame}{Q170: Ransomware Defenses}
		\begin{block}{Question}
			Which combination BEST defends against ransomware?
		\end{block}
		\begin{enumerate}[(A)]
			\item One antivirus tool
			\item Backups, segmentation, patching, email filtering, user training, and immutable copies
			\item Only MFA
			\item Only firewalls
		\end{enumerate}
	\end{frame}
	
	\begin{frame}{Q170: Explanation}
		\begin{block}{Correct Answer: B}
			\textbf{Ransomware defense is layered:}
			\begin{itemize}
				\item Backups for recovery
				\item Segmentation limits spread
				\item Patching closes entry points
				\item Training prevents initial access
			\end{itemize}
		\end{block}
		\begin{alertblock}{Critical Trap}
			Ransomware defense is layered, not a single tool.
		\end{alertblock}
	\end{frame}
	
	\begin{frame}{Q171: Patch vs Configuration}
		\begin{block}{Question}
			Which scenario is a PATCHING issue rather than a CONFIGURATION issue?
		\end{block}
		\begin{enumerate}[(A)]
			\item Firewall rule missing
			\item Known CVE unfixed on a server
			\item Weak password policy
			\item Shared admin account
		\end{enumerate}
	\end{frame}
	
	\begin{frame}{Q171: Explanation}
		\begin{block}{Correct Answer: B}
			\textbf{Unfixed CVE is a patch gap:}
			\begin{itemize}
				\item Known vulnerability remains
				\item Patch management addresses it
				\item Configuration issues are settings
			\end{itemize}
		\end{block}
		\begin{alertblock}{Critical Trap}
			CVEs are patch issues; settings are configuration issues.
		\end{alertblock}
	\end{frame}
	
	\begin{frame}{Q172: Vulnerability Prioritization}
		\begin{block}{Question}
			Which factor MOST improves vulnerability prioritization?
		\end{block}
		\begin{enumerate}[(A)]
			\item Alphabetical order
			\item Business criticality of the asset and exploit availability
			\item Age of the vulnerability
			\item Number of users}
		{B}
		\begin{itemize}
			\item
			\item Criticality drives impact
			\item Exploit availability drives likelihood
			\item Combine both
		\end{itemize}
		\begin{alertblock}{Critical Trap}
			Prioritize by business criticality and exploit availability.
		\end{alertblock}
	\end{frame}
	
	\begin{frame}{Q173: CVSS Limitations}
		\begin{block}{Question}
			What is a limitation of CVSS scoring?
		\end{block}
		\begin{enumerate}[(A)]
			\item It captures business context
			\item It does not capture business context, only technical severity
			\item It replaces prioritization
			\item It measures training
		\end{enumerate}
	\end{frame}
	
	\begin{frame}{Q173: Explanation}
		\begin{block}{Correct Answer: B}
			\textbf{CVSS is technical severity:}
			\begin{itemize}
				\item No asset criticality context
				\item Combine with business impact
				\item Use as one input
			\end{itemize}
		\end{block}
		\begin{alertblock}{Critical Trap}
			CVSS alone is incomplete. Add business context.
		\end{alertblock}
	\end{frame}
	
	\begin{frame}{Q174: Pen Test Rules of Engagement}
		\begin{block}{Question}
			What must rules of engagement define before a penetration test?
		\end{block}
		\begin{enumerate}[(A)]
			\item Only costs
			\item Scope, timing, escalation paths, and permitted techniques
			\item Only usernames
			\item Only vendor names
		\end{enumerate}
	\end{frame}
	
	\begin{frame}{Q174: Explanation}
		\begin{block}{Correct Answer: B}
			\textbf{RoE protects everyone:}
			\begin{itemize}
				\item Scope prevents collateral damage
				\item Timing protects peak hours
				\item Escalation enables safe handling
			\end{itemize}
		\end{block}
		\begin{alertblock}{Critical Trap}
			Rules of engagement prevent chaos during pen tests.
		\end{alertblock}
	\end{frame}
	
	\begin{frame}{Q175: SIEM Purpose}
		\begin{block}{Question}
			What is the primary purpose of a SIEM system?
		\end{block}
		\begin{enumerate}[(A)]
			\item Encrypt data
			\item Correlate security events and alert on threats
			\item Back up files
			\item Manage HR}
		{B}
		\begin{itemize}
			\item
			\item Aggregates logs
			\item Correlates events
			\item Alerts analysts
		\end{itemize}
		\begin{alertblock}{Critical Trap}
			SIEM correlates and alerts, not just stores logs.
		\end{alertblock}
	\end{frame}
	
	\begin{frame}{Q176: SOC Role}
		\begin{block}{Question}
			What does a Security Operations Center primarily do?
		\end{block}
		\begin{enumerate}[(A)]
			\item Marketing
			\item Monitor, detect, and respond to security incidents
			\item Manage HR
			\item Sell products}
		{B}
		\begin{itemize}
			\item
			\item Detect threats
			\item Analyze alerts
			\item Respond to incidents
		\end{itemize}
		\begin{alertblock}{Critical Trap}
			SOC is about monitoring and response.
		\end{alertblock}
	\end{frame}
	
	\begin{frame}{Q177: Threat Hunting vs Monitoring}
		\begin{block}{Question}
			What is threat hunting?
		\end{block}
		\begin{enumerate}[(A)]
			\item Waiting for alerts
			\item Proactively searching for hidden threats not caught by detection
			\item Buying tools
			\item Ignoring alerts}
		{B}
		\begin{itemize}
			\item
			\item Hypothesis-driven investigation
			\item Goes beyond alerts
			\item Finds stealthy attackers
		\end{itemize}
		\begin{alertblock}{Critical Trap}
			Hunting is proactive; monitoring is reactive.
		\end{alertblock}
	\end{frame}
	
	\begin{frame}{Q178: Threat Intelligence}
		\begin{block}{Question}
			What is threat intelligence used for?
		\end{block}
		\begin{enumerate}[(A)]
			\item Storing data
			\item Understanding adversary tactics, techniques, and indicators
			\item Measuring training
			\item Computing budgets}
		{B}
		\begin{itemize}
			\item
			\item Informs defenses
			\item Prioritizes patching
			\item Improves detection
		\end{itemize}
		\begin{alertblock}{Critical Trap}
			Threat intel guides defenses and prioritization.
		\end{alertblock}
	\end{frame}
	
	\begin{frame}{Q179: Incident Response Communications}
		\begin{block}{Question}
			Why is incident communications planning important?
		\end{block}
		\begin{enumerate}[(A)]
			\item Marketing
			\item Provides clear, timely information to stakeholders during an incident
			\item Reduces cost
			\item Avoids escalation}
		{B}
		\begin{itemize}
			\item
			\item Prevents confusion
			\item Supports trust
			\item Protects reputation
		\end{itemize}
		\begin{alertblock}{Critical Trap}
			Communication is part of IR, not an afterthought.
		\end{alertblock}
	\end{frame}
	
	\begin{frame}{Q180: Post-Incident Review}
		\begin{block}{Question}
			What is the primary purpose of a post-incident review?
		\end{block}
		\begin{enumerate}[(A)]
			\item Assign blame
			\item Identify lessons and improve future response
			\item Close tickets
			\item Reduce budgets}
		{B}
		\begin{itemize}
			\item
			\item Root causes
			\item Lessons learned
			\item Continuous improvement
		\end{itemize}
		\begin{alertblock}{Critical Trap}
			Post-incident reviews are for learning, not blame.
		\end{alertblock}
	\end{frame}
	
	\begin{frame}{Q181: Privacy Impact Assessment}
		\begin{block}{Question}
			When should a Privacy Impact Assessment be conducted?
		\end{block}
		\begin{enumerate}[(A)]
			\item After a breach
			\item Before implementing new processing that may affect privacy
			\item Only at year-end
			\item Never}
		{B}
		\begin{itemize}
			\item
			\item Early assessment
			\item Identifies privacy risks
			\item Shapes design
		\end{itemize}
		\begin{alertblock}{Critical Trap}
			PIAs are done before processing, not after.
		\end{alertblock}
	\end{frame}
	
	\begin{frame}{Q182: DPO Role}
		\begin{block}{Question}
			What is the primary role of a Data Protection Officer?
		\end{block}
		\begin{enumerate}[(A)]
			\item Marketing
			\item Advise on GDPR compliance and monitor compliance
			\item Manage IT
			\item Sell products}
		{B}
		\begin{itemize}
			\item
			\item Advises on compliance
			\item Monitors GDPR compliance
			\item Point of contact for regulators
		\end{itemize}
		\begin{alertblock}{Critical Trap}
			DPO advises and monitors, not operates.
		\end{alertblock}
	\end{frame}
	
	\begin{frame}{Q183: Cross-Border Data Flow}
		\begin{block}{Question}
			What is required for lawful cross-border data transfer from EU?
		\end{block}
		\begin{enumerate}[(A)]
			\item Nothing
			\item Valid legal mechanism such as adequacy decision or SCCs
			\item Only consent
			\item Email confirmation}
		{B}
		\begin{itemize}
			\item
			\item Adequacy decision
			\item Standard Contractual Clauses
			\item Binding Corporate Rules
		\end{itemize}
		\begin{alertblock}{Critical Trap}
			Cross-border transfers need legal mechanisms.
		\end{alertblock}
	\end{frame}
	
	\begin{frame}{Q184: Breach Notification to Individuals}
		\begin{block}{Question}
			When must individuals be notified of a breach under GDPR?
		\end{block}
		\begin{enumerate}[(A)]
			\item Always
			\item When the breach poses high risk to their rights and freedoms
			\item Never
			\item Only at year-end}
		{B}
		\begin{itemize}
			\item
			\item High-risk threshold
			\item Clear communication
			\item Support measures
		\end{itemize}
		\begin{alertblock}{Critical Trap}
			Individual notification depends on risk severity.
		\end{alertblock}
	\end{frame}
	
	\begin{frame}{Q185: HIPAA PHI Definition}
		\begin{block}{Question}
			What is PHI?
		\end{block}
		\begin{enumerate}[(A)]
			\item Public Health Information
			\item Protected Health Information, any information about health status linked to an individual
			\item Personal Hiring Information
			\item Private Internet}
		{B}
		\begin{itemize}
			\item
			\item Health information linked to individual
			\item Protected by HIPAA
			\item Broad definition
		\end{itemize}
		\begin{alertblock}{Critical Trap}
			PHI is health data plus identifier.
		\end{alertblock}
	\end{frame}
	
	\begin{frame}{Q186: PCI DSS Scope}
		\begin{block}{Question}
			What does PCI DSS cover?
		\end{block}
		\begin{enumerate}[(A)]
			\item Employee records
			\item Systems that store, process, or transmit cardholder data
			\item Marketing data
			\item Office supplies}
		{B}
		\begin{itemize}
			\item
			\item Cardholder data environment
			\item Reduces scope
			\item Applies to merchants and processors
		\end{itemize}
		\begin{alertblock}{Critical Trap}
			PCI DSS scope is cardholder data.
		\end{alertblock}
	\end{frame}
	
	\begin{frame}{Q187: Segmentation Purpose}
		\begin{block}{Question}
			Why segment networks?
		\end{block}
		\begin{enumerate}[(A)]
			\item Increase speed
			\item Limit lateral movement and reduce attack surface
			\item Save cost
			\item Reduce training}
		{B}
		\begin{itemize}
			\item
			\item Limits spread
			\item Isolates sensitive systems
			\item Supports compliance
		\end{itemize}
		\begin{alertblock}{Critical Trap}
			Segmentation limits blast radius.
		\end{alertblock}
	\end{frame}
	
	\begin{frame}{Q188: Data Classification Rules}
		\begin{block}{Question}
			Who is responsible for classifying data correctly?
		\end{block}
		\begin{enumerate}[(A)]
			\item The custodian
			\item The data owner or data steward
			\item The auditor
			\item The vendor
		\end{enumerate}
	\end{frame}
	
	\begin{frame}{Q188: Explanation}
		\begin{block}{Correct Answer: B}
			\textbf{Owners classify:}
			\begin{itemize}
				\item Data owners and stewards classify
				\item Custodians enforce
				\item Auditors verify
			\end{itemize}
		\end{block}
		\begin{alertblock}{Critical Trap}
			Owners classify; custodians enforce.
		\end{alertblock}
	\end{frame}
	
	\begin{frame}{Q189: Data Loss Prevention Modes}
		\begin{block}{Question}
			What are common DLP modes?
		\end{block}
		\begin{enumerate}[(A)]
			\item Fast and slow
			\item Monitor, block, and quarantine
			\item Old and new
			\item Red and green}
		{B}
		\begin{itemize}
			\item
			\item Monitor-only mode
			\item Blocking mode
			\item Quarantine mode
		\end{itemize}
		\begin{alertblock}{Critical Trap}
			DLP modes = monitor, block, quarantine.
		\end{alertblock}
	\end{frame}
	
	\begin{frame}{Q190: Key Rotation}
		\begin{block}{Question}
			Why is cryptographic key rotation important?
		\end{block}
		\begin{enumerate}[(A)]
			\item Saves money
			\item Limits exposure if a key is compromised
			\item Reduces storage
			\item Avoids backups}
		{B}
		\begin{itemize}
			\item
			\item Limits blast radius
			\item Complies with standards
			\item Better than never rotating
		\end{itemize}
		\begin{alertblock}{Critical Trap}
			Regular rotation limits key compromise damage.
		\end{alertblock}
	\end{frame}
	
	\begin{frame}{Q191: HSM vs KMS}
		\begin{block}{Question}
			What is the main difference between HSM and KMS?
		\end{block}
		\begin{enumerate}[(A)]
			\item They are the same
			\item HSM is dedicated hardware; KMS is software service managing keys
			\item KMS is hardware only
			\item HSM is software only}
		{B}
		\begin{itemize}
			\item
			\item HSM = dedicated hardware security
			\item KMS = key management service
			\item Both protect keys
		\end{itemize}
		\begin{alertblock}{Critical Trap}
			HSM is hardware; KMS is often software.
		\end{alertblock}
	\end{frame}
	
	\begin{frame}{Q192: SSL vs TLS}
		\begin{block}{Question}
			Which statement is correct about SSL and TLS?
		\end{block}
		\begin{enumerate}[(A)]
			\item They are the same
			\item TLS is the successor to SSL and is more secure
			\item SSL is newer than TLS
			\item TLS is deprecated}
		{B}
		\begin{itemize}
			\item
			\item TLS replaced SSL
			\item SSL is deprecated
			\item Use TLS versions
		\end{itemize}
		\begin{alertblock}{Critical Trap}
			TLS supersedes SSL.
		\end{alertblock}
	\end{frame}
	
	\begin{frame}{Q193: Certificate Pinning}
		\begin{block}{Question}
			What is certificate pinning?
		\end{block}
		\begin{enumerate}[(A)]
			\item Pinning a certificate to a bulletin board
			\item Binding a certificate to an application to prevent MITM attacks
			\item Deleting certificates
			\item Renewing certificates}
		{B}
		\begin{itemize}
			\item
			\item App trusts only specific certs
			\item Prevents MITM
			\item Common in mobile apps
		\end{itemize}
		\begin{alertblock}{Critical Trap}
			Pinning prevents unauthorized certificates.
		\end{alertblock}
	\end{frame}
	
	\begin{frame}{Q194: Zero-Day Vulnerability}
		\begin{block}{Question}
			What is a zero-day vulnerability?
		\end{block}
		\begin{enumerate}[(A)]
			\item A vulnerability patched yesterday
			\item A vulnerability unknown to the vendor with no patch available
			\item A training gap
			\item A configuration error}
		{B}
		\begin{itemize}
			\item
			\item Unknown to vendor
			\item No patch
			\item Requires compensating controls
		\end{itemize}
		\begin{alertblock}{Critical Trap}
			Zero-day = no patch yet. Compensate.
		\end{alertblock}
	\end{frame}
	
	\begin{frame}{Q195: Defense in Depth Example}
		\begin{block}{Question}
			Which is an example of defense in depth?
		\end{block}
		\begin{enumerate}[(A)]
			\item Only antivirus
			\item Firewalls, endpoint protection, MFA, backups, and training
			\item Only email filtering
			\item Only backups}
		{B}
		\begin{itemize}
			\item
			\item Layers cover different risks
			\item Reduces single-point failure
			\item Harder for attackers
		\end{itemize}
		\begin{alertblock}{Critical Trap}
			Defense in depth = multiple independent layers.
		\end{alertblock}
	\end{frame}
	
	\begin{frame}{Q196: Least Privilege Example}
		\begin{block}{Question}
			Which is a real-world example of least privilege?
		\end{block}
		\begin{enumerate}[(A)]
			\item Giving all staff admin rights
			\item Giving support staff read-only access unless escalation is approved
			\item Sharing admin passwords
			\item Allowing everyone to approve payments}
		{B}
		\begin{itemize}
			\item
			\item Minimum access required
			\item Escalation for more access
			\item Prevents abuse
		\end{itemize}
		\begin{alertblock}{Critical Trap}
			Least privilege is minimum access, not maximum.
		\end{alertblock}
	\end{frame}
	
	\begin{frame}{Q197: Access Review Frequency}
		\begin{block}{Question}
			How often should access reviews be conducted for critical systems?
		\end{block}
		\begin{enumerate}[(A)]
			\item Every five years
			\item At least annually, and more frequently for critical systems
			\item Only when someone leaves
			\item Never}
		{B}
		\begin{itemize}
			\item
			\item Annual minimum
			\item More often for critical
			\item After major changes
		\end{itemize}
		\begin{alertblock}{Critical Trap}
			Critical systems need more frequent access reviews.
		\end{alertblock}
	\end{frame}
	
	\begin{frame}{Q198: Encryption Key Escrow}
		\begin{block}{Question}
			What is key escrow?
		\end{block}
		\begin{enumerate}[(A)]
			\item Sending keys by email
			\item Storing keys with a trusted third party for recovery
			\item Deleting keys
			\item Sharing keys publicly}
		{B}
		\begin{itemize}
			\item
			\item Enables recovery
			\item Used for encryption keys
			\item Requires strong controls
		\end{itemize}
		\begin{alertblock}{Critical Trap}
			Key escrow requires trust and controls.
		\end{alertblock}
	\end{frame}
	
	\begin{frame}{Q199: Privacy by Default}
		\begin{block}{Question}
			What does privacy by default mean?
		\end{block}
		\begin{enumerate}[(A)]
			\item All data is public
			\item Systems default to the most privacy-protective settings
			\item Users must opt out of sharing
			\item Ignore privacy}
		{B}
		\begin{itemize}
			\item
			\item Privacy-protective defaults
			\item Users opt in
			\item Common in GDPR
		\end{itemize}
		\begin{alertblock}{Critical Trap}
			Privacy by default = strictest settings by default.
		\end{alertblock}
	\end{frame}
	
	\begin{frame}{Q200: Ultimate Security Goal}
		\begin{block}{Question}
			What is the ultimate goal of information security?
		\end{block}
		\begin{enumerate}[(A)]
			\item Eliminate all attacks
			\item Enable business objectives while managing information risk to acceptable levels
			\item Maximize tools
			\item Increase complexity}
		{B}
		\begin{itemize}
			\item
			\item Security enables business
			\item Risk is managed to acceptable levels
			\item Absolute security is impossible
		\end{itemize}
		\begin{alertblock}{Critical Trap}
			Security enables business; it is not an end in itself.
		\end{alertblock}
	\end{frame}
	
	\begin{frame}{Exam Mindset Summary}
		\begin{block}{Remember}
			\begin{itemize}
				\item \textbf{CIA triad:} Confidentiality, Integrity, Availability
				\item \textbf{IAAA:} Identification, Authentication, Authorization, Accountability
				\item \textbf{Encryption is reversible; hashing is one-way}
				\item \textbf{Signatures} provide integrity, authenticity, non-repudiation -- NOT confidentiality
				\item \textbf{Privacy} is about rights; \textbf{security} is about protection
				\item \textbf{Zero Trust:} verify every access, no implicit trust
				\item \textbf{Defense in depth:} multiple independent layers
				\item \textbf{Cloud:} shared responsibility; data location matters legally
				\item \textbf{DLP, PAM, SIEM, WAF, HSM:} know what each does and does NOT do
			\end{itemize}
		\end{block}
	\end{frame}
	
	\begin{frame}{Thank You}
		\begin{center}
			{\Huge \textbf{Good Luck on Your CRISC Exam!}}
		\end{center}
	\end{frame}
	
\end{document}